\section{Background}
\label{sec:bg}

AIMer is a post-quantum digital signature scheme. Its security relies on a symmetric-key one-way function. This one-way function is called AIM2 and is defined over the binary field $\gftwo{\lambda}$. A one-way function is easy to evaluate, but it is difficult to recover its input from its output. In this sense, it is similar to a hash function.

The secret key is a randomly chosen input value. The public key is the output obtained by applying AIM2 to this input. A signature proves that the signer knows the input value that generated the public key. At the same time, it does not reveal the secret input itself. This type of proof, which proves only the fact of knowledge without revealing the secret, is called a \textbf{zero-knowledge proof}. AIMer generates signatures using a zero-knowledge proof~\cite{fiege1987zero}.

AIMer constructs this proof using the \textbf{MPC-in-the-Head (MPCitH)} technique~\cite{delpech2021limbo}. It simulates a setting in which $N$ virtual parties hold shares of the secret key and jointly compute AIM2. The verifier then checks this simulated computation. A single execution still leaves room for cheating. Therefore, the same process is repeated several dozen times, denoted by $\tau$ in Table~\ref{tab:params}.

Two types of operations dominate the signing cost. First, the $N$ parties repeatedly perform \emph{the same binary-field arithmetic}. Second, \emph{Keccak hashing}~\cite{bertoni2009keccak,bertoni2013keccak} is used extensively throughout the process for commitments and challenge generation. This paper accelerates these two components using VPCLMULQDQ and AVX-512VL (Vector Length)~\cite{anderson2018enhanced}, respectively. The parameters of each variant are shown in Table~\ref{tab:params}.

\begin{table}[h]
\centering
\caption{Parameters of AIMer variants. $\lambda$: security bits/field size, $N$: number of MPC parties,
$\tau$: number of parallel repetitions.}
\label{tab:params}
{\normalsize
\setlength{\tabcolsep}{7pt} % column padding (default is 6pt)
\renewcommand{\arraystretch}{1.15} % row height
\begin{tabular}{lrrrrr}
\toprule
Variant & Security level & Field $\gftwo{\lambda}$ & $N$ & $\tau$ & Signature (bytes) \\
\midrule
AIMER-128f & NIST~1 & $\gftwo{128}$ & 16  & 33 & 5{,}888 \\
AIMER-128s & NIST~1 & $\gftwo{128}$ & 256 & 17 & 4{,}160 \\
AIMER-192f & NIST~3 & $\gftwo{192}$ & 16  & 49 & 13{,}056 \\
AIMER-192s & NIST~3 & $\gftwo{192}$ & 256 & 25 & 9{,}120 \\
AIMER-256f & NIST~5 & $\gftwo{256}$ & 16  & 65 & 25{,}120 \\
AIMER-256s & NIST~5 & $\gftwo{256}$ & 256 & 33 & 17{,}056 \\
\bottomrule
\end{tabular}
}
\end{table}

\section{Related Work}
\label{relatedwork}

Prior work related to this paper can be divided into two lines. The first line is \emph{MPCitH-based signatures}, to which AIMer belongs. This line shows how signatures have been made smaller and faster by choosing suitable symmetric primitives and checking techniques~\cite{lee2023high}. The second line is SIMD-based cryptographic optimization, which this paper applies to AIMer. This line focuses on accelerating cryptographic operations using vector instructions. We review these two lines in order.

\subsection{MPCitH-Based Signatures}

MPC-in-the-Head (MPCitH) is a zero-knowledge proof paradigm proposed by Ishai et al.~\cite{ishai2007zero}. In this paradigm, the proof size is proportional to the number of \emph{multiplications}, or nonlinear gates, in the circuit. Subsequent studies have reduced this cost in several directions.

First, Picnic adopted LowMC as a one-way function. LowMC is a block cipher with a small number of AND gates, which reduces the number of multiplications itself~\cite{chase2017post}. Second, later work introduced techniques that trade communication cost and signature size by using a preprocessing model and adjusting the number of parties $N$. These techniques reduced signatures to the range of several tens of kilobytes~\cite{katz2018improved}. Third, Banquet used the standard block cipher AES instead of LowMC. Instead of verifying inverse-S-box multiplications one by one, it grouped them into polynomial identities and checked them at once using an amortized check. This reduced the cost of multiplication verification~\cite{baum2021banquet}. BN++ further simplified this checking procedure~\cite{kales2022efficient}. This last approach was adopted by AIMer.

Furthermore, the FAEST digital signature scheme uses VOLE-in-the-Head~\cite{baum2023publicly} instead of MPCitH. It efficiently generates correlated randomness and reduces AES-based signatures to about 5~KB. FAEST advanced to the second round of the NIST additional signature competition~\cite{baum2025faest}. In this line of work, AIMer is a case that introduces a \emph{dedicated} one-way function, AIM. AIM obtains short signatures by minimizing the number of multiplications in the circuit with a low-multiplication S-box. After AIM was published, it was subjected to algebraic analysis~\cite{liu2023algebraic}. AIM2, a strengthened version reflecting this analysis, was subsequently adopted in AIMer, which was selected as a final algorithm in the KpqC competition~\cite{choaimer}. Another line that obtains security only from symmetric primitives is the hash-based signature family SPHINCS+/SLH-DSA~\cite{bernstein2019sphincs+,saarinen2024accelerating}.

\subsection{SIMD-Based Cryptographic Optimization}

SIMD acceleration has been actively studied for compute-intensive
workloads such as sorting~\cite{chhugani2008efficient}. In cryptographic software, Drucker et al. used the fact that 512-bit VPCLMULQDQ performs four independent $64\times64$ carry-less multiplications in a single instruction. With this property, they substantially improved the throughput of binary-field multiplication and GHASH compared with PCLMULQDQ~\cite{drucker2018fast}. The same authors also proposed VAES. VAES vectorizes AES rounds by packing four blocks into one register. This accelerates parallel modes such as AES-GCM~\cite{drucker2019making}.

SIMD optimization has also been studied for lattice-based PQC. For example, AVX-512 vectorization of Dilithium has been reported~\cite{xu2024optimizing}. For Korean standard cryptographic algorithms, AVX-512 acceleration has been reported for the LSH hash function and LEA-128-GCM. LEA-GCM combines 16-block parallel CTR mode with four-block VPCLMULQDQ-based GHASH. It achieved about a $3.3\times$ improvement over the KISA SSE reference implementation. LSH achieved an average speedup of $2.28\times$ using dual-message processing based on 512-bit register interleaving~\cite{lee2026advanced}.

Prior implementation studies have also examined AIM and AIM2 directly. Lee et al. accelerated the original AIM symmetric primitive by optimizing the Mer operation and simplifying its linear layer~\cite{lee2023high}. Song et al. designed quantum circuits for all AIM2 variants and analyzed their resource requirements~\cite{song2025quantum}. These studies respectively address software optimization of the underlying AIM primitive and quantum resource analysis, but neither presents an AVX-512 implementation covering all six AIMer variants. A recent SIMD study in a related proof-system setting accelerated FAEST on ARM using NEON AES instructions and parallel VOLE generation~\cite{lee2026accelerating}. In contrast, our work provides an end-to-end AVX-512 implementation of all six AIMer variants, jointly optimizing binary-field arithmetic, the linear layer, and Keccak.